\documentclass{article}
\usepackage[utf8]{inputenc}
\usepackage[brazil]{babel}
\usepackage{tikz}
\usepackage{amssymb}
\usepackage{helvet}
\renewcommand{\familydefault}{\sfdefault}
\pagestyle{empty}
\begin{document}
\begin{tikzpicture}[scale=0.65, every node/.style={font=\scriptsize}]
    % Padding da sync_matrix (linha 0 e coluna 0)
    \fill[gray!20] (0,4) rectangle (5,5);
    \fill[gray!20] (0,0) rectangle (1,4);
    \node[gray!80!black] at (2.5, 4.5) {\textit{padding} (\texttt{sync\_matrix})};

    % Blocos ja concluidos
    \fill[green!20] (1,3) rectangle (2,4);
    \fill[green!20] (2,3) rectangle (3,4);
    \fill[green!20] (1,2) rectangle (2,3);

    % Frente de onda ativa (anti-diagonal em paralelo)
    \fill[blue!35] (3,3) rectangle (4,4);
    \fill[blue!35] (2,2) rectangle (3,3);
    \fill[blue!35] (1,1) rectangle (2,2);

    % Blocos futuros (aguardando dependencias)
    \fill[white] (4,3) rectangle (5,4);
    \fill[white] (3,2) rectangle (5,3);
    \fill[white] (2,1) rectangle (5,2);
    \fill[white] (1,0) rectangle (5,1);

    % Grade e contornos
    \draw[step=1, gray] (0,0) grid (5,5);
    \draw[thick] (1,0) rectangle (5,4);

    % Rotulos dos blocos
    \node at (1.5,3.5) {$\checkmark$};
    \node at (2.5,3.5) {$\checkmark$};
    \node at (1.5,2.5) {$\checkmark$};
    \node at (3.5,3.5) {\textbf{T1}};
    \node at (2.5,2.5) {\textbf{T2}};
    \node at (1.5,1.5) {\textbf{T3}};

    % Setas de dependencia para o bloco T2 (2,2)
    \draw[->, thick, red!75!black] (2.5, 3.2) -- (2.5, 2.8);
    \draw[->, thick, red!75!black] (1.8, 2.5) -- (2.2, 2.5);

    % Legenda lateral
    \fill[green!20] (5.6, 3.6) rectangle (6.1, 4.0);
    \draw[gray] (5.6, 3.6) rectangle (6.1, 4.0);
    \node[anchor=west] at (6.2, 3.8) {Blocos concluídos};

    \fill[blue!35] (5.6, 2.9) rectangle (6.1, 3.3);
    \draw[gray] (5.6, 2.9) rectangle (6.1, 3.3);
    \node[anchor=west] at (6.2, 3.1) {Frente diagonal ativa (paralelo)};

    \draw[->, thick, red!75!black] (5.6, 2.4) -- (6.1, 2.4);
    \node[anchor=west] at (6.2, 2.4) {Dependências \texttt{in/out}};

    \node[anchor=west, align=left] at (5.6, 1.3) {Dentro de cada bloco:\\cálculo sequencial por linha\\(localidade espacial em cache)};
\end{tikzpicture}
\end{document}
